\subsection{单代数的可对角化子代数}

设 D 是一个是体的 K-代数，V 是 D 上一个有限维右向量空间。设 L 是 \(\operatorname{End}_{\mathrm{D}}(\mathrm{V})\) 的一个子代数且是一个可对角化 K-代数（V, §6, No.~3, 第 28 页）。按定义，L 在 K 上是有限次数的，且存在 L 在 K 上的一个基 \((\varepsilon_{i})_{i\in\mathrm{I}}\)，它具有下列性质：

\[
\varepsilon_ {i} ^ {2} = \varepsilon_ {i}, \varepsilon_ {i} \varepsilon_ {j} = 0 \text { if } i \neq j, \quad \sum_ {i \in I} \varepsilon_ {i} = 1.
\]

对 I 中的每个 i，设 \(V_{i} = \varepsilon_{i}(V)\)；则 \((\mathrm{V}_{i})_{i \in \mathrm{I}}\) 是 V 的一族非零线性子空间，其直和为 V（II, §1, No.~8, 第 209 页，命题 12）。设 u 是 V 的一个自同态；则 u 属于 L 当且仅当对每个 \(i \in I\)，存在 K 的一个元素 \(\lambda_{i}\)，使得 \(u(x) = \lambda_{i}x\)（对每个 \(x \in V_{i}\)）。

反之，假设 V 是一族非零线性子空间 \((\mathrm{V}_{i})_{i\in\mathrm{I}}\) 的直和。对任何元素 \(\boldsymbol{\lambda}=(\lambda_{i})_{i\in\mathrm{I}}\)（\(K^{I}\) 的），用 \(u_{\lambda}\) 表示 D-向量空间 V 的这样的自同态：\(u_{\boldsymbol{\lambda}}(x)=\lambda_{i}x\)（对 \(x\in V_{i}\)）。诸自同态 \(u_{\lambda}\)（\(\lambda\in K^{I}\)）所成的集 L 是 \(\operatorname{End}_{\mathrm{D}}(\mathrm{V})\) 的一个可对角化子代数，它以族 \((\varepsilon_{i})_{i\in\mathrm{I}}\) 为基，其中 \(\varepsilon_{i}\) 是以 \(V_{i}\) 为像、以 \(\sum_{j\neq i}V_{j}\) 为核的射影。我们称 L 为 \(\operatorname{End}_{\mathrm{D}}(\mathrm{V})\) 中与直和分解 \(V=\oplus_{i\in\mathrm{I}}V_{i}\) 相伴的可对角化子代数。我们有 \([L:K]=\operatorname{Card}(I)\leqslant\dim_{\mathrm{D}}(V)\)。

命题 6. —— 设 L 是 \(\operatorname{End}_{\mathrm{D}}(\mathrm{V})\) 的可对角化子代数，它与直和分解 \(V = \oplus_{i \in I} V_i\) 相伴。

\begin{enumerate}
\def\labelenumi{\alph{enumi})}
\item
  代数 L 在 K-代数 \(\operatorname{End}_{\mathrm{D}}(\mathrm{V})\) 的可对角化子代数中是极大的当且仅当每个 \(V_{i}\) 在 D 上的维数为 1。
\item
  代数 \(\mathrm{L}\) 是 \(\operatorname{End}_{\mathrm{D}}(\mathrm{V})\) 的极大交换子代数当且仅当我们有 \(\mathrm{D} = \mathrm{K}\) 且每个 \(\mathrm{V}_i\) 在 \(\mathrm{K}\) 上的维数为 1。
\end{enumerate}

若每个向量空间 \(V_{i}\) 在 D 上的维数为 1，则我们有

\[
[ \mathrm{L}: \mathrm{K} ] = \operatorname{Card} (\mathrm{I}) = \dim_ {\mathrm{D}} (\mathrm{V}),
\]

于是 L 在 \(\operatorname{End}_{\mathrm{D}}(\mathrm{V})\) 的可对角化子代数中是极大的。在相反的情形，存在一个指标 \(j \in I\)，使得 \(\dim_{\mathrm{D}}(\mathrm{V}_{j}) \geqslant 2\)。选取两个非零线性子空间 \(V_{j}^{\prime}\) 与 \(V_{j}^{\prime\prime}\)（\(V_{j}\) 的），使其直和为 \(V_{j}\)。\(\operatorname{End}_{\mathrm{D}}(\mathrm{V})\) 中与直和分解 \(\mathrm{V} = (\oplus_{i \in \mathbf{I} - \{j\}} \mathrm{V}_{i}) \oplus \mathrm{V}_{j}^{\prime} \oplus \mathrm{V}_{j}^{\prime\prime}\) 相伴的可对角化子代数包含 L 且不等于 L；由此得到断言 a)。

交换子 \(L'\)（L 在 \(\operatorname{End}_{\mathrm{D}}(\mathrm{V})\) 中的）由形如 \((x_{i}) \mapsto (u_{i}(x_{i}))\) 的自同态组成，其中 \((u_{i}) \in \prod_{i \in \mathrm{I}} \operatorname{End}_{\mathrm{D}}(\mathrm{V}_{i})\)。代数 L 是 \(\operatorname{End}_{\mathrm{D}}(\mathrm{V})\) 的极大交换子代数当且仅当我们有 \(L = L'\)（VIII, 第 261 页，引理 3, a)）。于是这一关系等价于“\(\operatorname{End}_{\mathrm{D}}(\mathrm{V}_{i}) = \mathrm{K}\)（对每个 \(i \in I''\)）”；由此得到断言 b)。

命题 7. —— 设 L 是 K 上一个有限次数的交换代数。下列性质等价：

\begin{enumerate}
\def\labelenumi{(\roman{enumi})}
\item
  代数 L 是平展的。
\item
  存在 K 的一个有限次数的可分扩张，它将 K 对角化。
\end{enumerate}

蕴含关系 \((\mathrm{ii})\Rightarrow(\mathrm{i})\) 由 V, §6, No.~3, 第 29 页的命题 2 得出。我们来证明蕴含关系 \((\mathrm{i})\Rightarrow(\mathrm{ii})\)。设 \(\Omega\) 是 K 的一个可分闭包。由 V, §6, No.~7, 第 34 页的定理 4，存在 K 的有限次数扩张 \(L_{1},\ldots,L_{n}\)，它们包含于 \(\Omega\)，使得 L 同构于积

\(L_{1} \times \cdots \times L_{n}\)。设 N 是 K 的一个包含诸 \(L_{i}\) 的伽罗瓦扩张（V, §10, No.~1, 第 58 页），我们来证明 \(L_{(N)}\) 是可对角化的。由本原元定理（V, §7, No.~4, 第 40 页，定理 1），对每个 \(i \in [1, n]\)，存在一个不可约的可分多项式 \(P_{i} \in K[X]\)，使得 \(L_{i}\) 同构于 \(K[X]/(P_{i})\)。由于 N 是 K 的一个正规扩张且 \(P_{i}\) 在其中有一个根，多项式 \(P_{i}\) 在 N{[}X{]} 中分裂，且根为单根。于是，N-代数 \(L_{i(N)}\)（它同构于 \(N[X]/(P_{i})\)）同构于 \(N^{[L_{i}:K]}\)。因此 \(L_{(N)}\) 是可对角化的。

定理 7. —— 设 A 是一个有限次数的中心单 K-代数，L 是 A 的一个子代数。下列性质等价：

\begin{enumerate}
\def\labelenumi{(\roman{enumi})}
\item
  代数 L 是 A 的一个极大平展子代数。
\item
  存在 K 的一个扩张 \(K'\)、一个整数 \(n \geqslant 1\)，以及一个同构 \(\theta\)（从 \(\mathrm{A}_{(\mathrm{K}')}\) 到 \(\mathbf{M}_{n}(\mathrm{K}')\) 的），它把 \(\mathrm{L}_{(\mathrm{K}')}\) 映到对角矩阵所成之集。
\item
  存在如 (ii) 中的 \(K'\)、n 与 \(\theta\)，其中扩张 \(K'\) 还被假设是伽罗瓦的且是有限次数的。
\end{enumerate}

显然 (iii) 蕴含 (ii)。

若性质 (ii) 成立，则 \(L_{(K')}\) 是 \(A_{(K')}\) 的一个极大交换子代数（命题 6），且它是可对角化的。于是 K-代数 L 是平展的（V, §6, No.~3, 第 28 页，定义 1），且它是 A 的一个极大交换子代数（VIII, 第 261 页，引理 3, b)）。我们证明了 (ii) 蕴含 (i)。

假设性质 (i) 成立。由于 L 在 K 上是平展的，由命题 7，存在 K 的一个有限次数的伽罗瓦扩张 \(K_{1}\)，使得 \(K_{1}\)-代数 \(\mathrm{L}_{(\mathrm{K}_{1})}\) 是可对角化的。代数 A 是中心单的；由（VIII, 第 252 页，定理 1），存在一个伽罗瓦扩张 \(K_{2}\)、一个 \(K_{2}\) 上维数为 n 的向量空间 V，以及一个同构 \(\theta\)（从 \(\mathrm{A}_{(\mathrm{K}_{2})}\) 到 \(\mathrm{End}_{\mathrm{K}_{2}}(\mathrm{V})\) 的）。由 V, 第 55 页的命题 1，我们可以假设 \(K_{1}=K_{2}\)。由 VIII, 第 264 页的命题 4, b) 与 VIII, 第 261 页的引理 3, b)，\(\mathrm{L}_{(\mathrm{K}^{\prime})}\) 是 \(\mathrm{A}_{(\mathrm{K}^{\prime})}\) 的一个极大交换子代数，于是 \(\theta(\mathrm{L}_{(\mathrm{K}^{\prime})})\) 是 \(\mathrm{End}_{\mathrm{K}^{\prime}}(\mathrm{V}^{\prime})\) 的一个极大交换子代数。我们对可对角化代数 \(\theta(\mathrm{L}_{(\mathrm{K}^{\prime})})\) 应用命题 6：存在一个基 \((e_{1},\ldots,e_{n})\)（\(V^{\prime}\) 在 \(K^{\prime}\) 上的），使得 \(\theta(\mathrm{L}_{(\mathrm{K}^{\prime})})\) 由 \(V^{\prime}\) 的关于该基具有对角矩阵的自同态组成。因此 (i) 蕴含 (iii)。

